% Generated draft: requires mathematical and source review.
\documentclass[11pt]{article}
\usepackage[margin=25mm]{geometry}
\usepackage{fontspec}
\setmainfont{Latin Modern Roman}
\usepackage{amsmath,amssymb,mathtools,unicode-math}
\setmathfont{Latin Modern Math}
\usepackage{enumitem,needspace,xurl,hyperref}
\setlength{\parindent}{0pt}
\setlength{\parskip}{5pt}
\setlength{\emergencystretch}{3em}
\newcommand{\N}{\mathbb{N}}
\newcommand{\Z}{\mathbb{Z}}
\newcommand{\Q}{\mathbb{Q}}
\newcommand{\R}{\mathbb{R}}
\newcommand{\C}{\mathbb{C}}
\newcommand{\F}{\mathbb{F}}
\newcommand{\A}{\mathbb{A}}
\newcommand{\PP}{\mathbb{P}}
\newcommand{\E}{\mathbb{E}}
\newcommand{\eps}{\varepsilon}
\newcommand{\dd}{\,\mathrm{d}}
\begin{document}
\section*{UnsolvedMath ID: 800389}
\textbf{Source problem code:} catalogue entry
\subsection{Definitions and mathematical statement}
Let \(K\) be an algebraically closed field of characteristic \(0\), let \(X\) be a quasi-projective variety over \(K\), and let \(f:X\to X\) be a \(K\)-endomorphism, meaning a morphism from \(X\) to itself defined over \(K\). For \(P\in X(K)\), define its forward orbit by
\[
\mathcal O_f(P)=\{f^n(P):n\in\Z_{\geq 0}\},\qquad f^0=\operatorname{id}_X.
\]
A subset of \(X(K)\) is Zariski dense in \(X\) if its Zariski closure is \(X\). Put
\[
D_f=\{P\in X(K):\overline{\mathcal O_f(P)}^{\mathrm{Zar}}=X\}.
\]
Assume, as in the supplied source, that \(f\) is ``not a fibered map.'' The supplied record does not define this term. The following assertions therefore retain that hypothesis with its intended source meaning; resolving its definition is necessary for a fully specified formulation. No dominance hypothesis on \(f\) is explicitly supplied.

\begin{enumerate}
\item[(a)] \textbf{Weak form.} For every \(K,X,f\) satisfying these hypotheses, there exists \(P\in X(K)\) such that
\[
\overline{\mathcal O_f(P)}^{\mathrm{Zar}}=X.
\]
\item[(b)] \textbf{Strong form.} For every \(K,X,f\) satisfying these hypotheses,
\[
\overline{D_f}^{\mathrm{Zar}}=X.
\]
\end{enumerate}
These are the two conjectural assertions recorded by the source, subject to the unresolved meaning of ``fibered map.''
\subsection{Short English statement}
For a self-morphism of a quasi-projective variety over an algebraically closed field of characteristic zero, assume the source's unspecified condition that the map is not fibered. The weak conjecture predicts a point with Zariski dense forward orbit; the strong conjecture predicts that such points are Zariski dense. The meaning of the nonfibered hypothesis requires clarification.
\subsection{Sources}
\begin{itemize}
\item Local UnsolvedMath registry: permanent ID \(800389\), source problem code \texttt{catalogue entry}, title ``Open Problems in Arithmetic Dynamics --- Conjecture 12.'' The supplied registry labels the problem open; this status is supplied metadata, not an independently verified mathematical or literature assessment.
\item Benedetto, DeMarco, Ingram, Jones et al., \emph{Current Trends and Open Problems in Arithmetic Dynamics} (2018), Conjecture 12. Supplied source URL: \url{https://arxiv.org/abs/1806.04980}. The background records access on 2026-10-16 and extraction by \texttt{arxiv\_latex}; these are supplied provenance statements, not a claim of independent access or verification.
\item The statement attributes the conjecture to Amerik--Bogomolov--Rovinsky (2011), reference label \texttt{MR2862064}, and Medvedev--Scanlon (2009), reference label \texttt{arxiv0901.2352}. Full bibliographic entries for these references are not supplied.
\end{itemize}
\end{document}
